% BEGIN R28_MEMORY_PARITY
% Continuación aditiva de MEM01: los lectores ya están definidos.
\begin{proposition}[Parité des lecteurs sous inversion orientée]
Avec des indices cycliques dans \(\mathbb Z/12\mathbb Z\), soient
\((C_M\tau)_m=-\tau_{-m}\) et \(\operatorname{sgn}(0)=0\).
Les deux lecteurs satisfont
\[
 \nu_{120}(C_M\tau)=-\nu_{120}(\tau),\qquad
 \nu_{270}(C_M\tau)=\nu_{270}(\tau).
\]
\end{proposition}
\begin{proof}
La réflexion \(m\mapsto-m\) permute les quatre classes de positions
modulo quatre. L'inversion change le signe de chaque somme de classe et,
par l'imparité de \(\operatorname{sgn}\), celui de \(\nu_{120}\).
Dans le second lecteur, les signes des deux facteurs se compensent :
\[
 \sum_m(C_M\tau)_m(C_M\tau)_{m+3}
 =\sum_m\tau_{-m}\tau_{-m-3}
 =\sum_k\tau_{k+3}\tau_k
 =\nu_{270}(\tau).
\]
\end{proof}
L'orientation distingue ainsi une lecture impaire et une corrélation paire.
Toutes deux dépendent de la disposition des échantillons dans le registre ; leurs
lois de transformation précisent quelle information persiste lorsqu'on inverse
la route.
% END R28_MEMORY_PARITY
